% fonts-takao-mincho (00303.01-2) 

% Janusz S. Bień
% Formal Linguistics Department, University of Warsaw, Dobra 55, 00-312 Warszawa, Poland
% jsbien@uw.edu.pl

\documentclass{mwart}
\usepackage{fontspec}
\usepackage{polyglossia}
\usepackage{etoolbox}
\setmainlanguage{polish}
\setotherlanguage{english}
\usepackage{csquotes}
\usepackage{wrapfig}
\usepackage[style=authoryear,backend=biber,backref=true,urldate=short]{biblatex}
%\addbibresource{jsbqed.bib}
% \addbibresource{4JSB2014u.bib}
% \addbibresource{2017.bib}
% \addbibresource{2018.bib}
% \addbibresource{2019a.bib}
% \addbibresource{2019b.bib}
\addbibresource{all.bib}
\addbibresource{tmp.bib}

\DefineBibliographyStrings{polish}{%
  byeditor = {red.},
  urlseen = {dostęp},
}

% http://tex.stackexchange.com/questions/166337/quotation-mark-quotation-sign-xelatex-polyglossia-csquotes
\DeclareQuoteStyle{polish}% I looked it up on Wikipedia, no idea if it's right
  {\quotedblbase}
  {\textquotedblright}
  [0.05em]
  {\textquoteleft}
  {\textquoteright}

% psuje tytuły???:
%  \usepackage{ulem}
\usepackage{metalogo}
\usepackage[polish]{varioref}

\def\eob{ę}

\setmainfont[Mapping=tex-text]{TeX Gyre Termes}
% \char"EC10 \char"EC11 oraz \char"EC12
\setmonofont{NotoSansMono-Regular}

\def\orogate{\char"EC12}
\usepackage{relsize}

\usepackage{hyperref}

\usepackage{graphicx}
% [hyphens]: options clash
\usepackage{url}
%\usepackage{natbib}
\newcommand{\uname}[1]{{\relsize{-1}\texttt{#1}}}
\newcommand{\ucode}[1]{\texttt{#1}}
\newcommand{\mname}[1]{\texttt{#1}}
\newcommand{\mcode}[1]{\texttt{#1}}
\newcommand{\aname}[1]{\texttt{#1}}
\newcommand{\acode}[1]{\texttt{#1}}
\newcommand{\prname}[1]{\textsf{#1}}
%


% \usepackage{draftwatermark}
\usepackage[doublespacing]{setspace}

\usepackage[xetex]{preview}

\begin{document}
%\gappto\captionslingua{\renewcommand{\chaptername}{Caput}}
\gappto\captionspolish{\renewcommand{\figurename}{Ilustracja}}

\renewcommand{\topfraction}{0.9}	
\setcounter{topnumber}{4}
\setcounter{totalnumber}{4}

% 
\title{Standard Unicode i dawne teksty polskie\\\relsize{-1}Uaktualniony postprint\footnote{Tekst ten
    stanowi autorską wersję artykułu, ktory został opublikowany bez
    korekty autorskiej z drobnymi zmianami i pod zmienionym tytułem w
    \textit{Acta Poligraphica} (ISSN 2299-9981) w numerze 14 z 2019~r.
    (strony 7-28). W artykule zostały ponownie zweryfikowane adresy
    internetowe i wprowadzono drobne uaktualnienia oraz
    uzupełnienia.}\\\relsize{-3}DOI:10.5281/zenodo.4058702}

\author{Janusz S. Bień}

%\date{\today}
\date{30 września 2020 r.}

\maketitle


\input JSB_tekstele-tekstPP

% redundantne:
\end{document}

%%% Local Variables: 
%%% coding: utf-8-unix
%%% mode: latex
%%% TeX-master: t
%%% TeX-PDF-mode: t
%%% TeX-engine: xetex
%%% End: 
